% Standalone problem document, generated from the adjacent README.md.
% Compile with XeLaTeX. Mathematical content is maintained in Markdown.
\documentclass[11pt,a4paper]{article}
\usepackage[fontsize=10.5pt]{scrextend}
\usepackage[margin=22mm,headheight=15pt,headsep=9mm,footskip=11mm]{geometry}
\usepackage{fontspec}
\setmainfont{texgyrepagella}[Extension=.otf,UprightFont=*-regular,BoldFont=*-bold,ItalicFont=*-italic,BoldItalicFont=*-bolditalic]
\setsansfont{texgyreheros}[Extension=.otf,UprightFont=*-regular,BoldFont=*-bold,ItalicFont=*-italic,BoldItalicFont=*-bolditalic]
\setmonofont{texgyrecursor}[Extension=.otf,UprightFont=*-regular,BoldFont=*-bold,ItalicFont=*-italic,BoldItalicFont=*-bolditalic,Scale=MatchLowercase]
\usepackage{amsmath,amssymb}
\usepackage{unicode-math}
\setmathfont{texgyrepagella-math.otf}
\usepackage{xcolor,microtype,enumitem,needspace,fancyhdr}
\usepackage{bookmark}
\definecolor{ink}{HTML}{17344B}
\definecolor{accent}{HTML}{197D80}
\definecolor{muted}{HTML}{596773}
\hypersetup{colorlinks=true,linkcolor=accent,urlcolor=accent,pdftitle={SP-15 - Finitely
many unitary classes with prescribed shifted singular
values},pdfauthor={Open Problems in Numerical Linear Algebra}}
\urlstyle{same}
\setlength{\parindent}{0pt}
\setlength{\parskip}{5pt plus 1pt minus 1pt}
\linespread{1.04}
\setlength{\emergencystretch}{3em}
\widowpenalty=10000
\clubpenalty=10000
\displaywidowpenalty=10000
\allowdisplaybreaks[1]
\setlist{nosep,leftmargin=1.5em}
\providecommand{\tightlist}{\setlength{\itemsep}{0pt}\setlength{\parskip}{0pt}}
\providecommand{\pandocbounded}[1]{#1}
\pagestyle{fancy}
\fancyhf{}
\fancyhead[L]{\sffamily\footnotesize\color{muted}OPEN PROBLEMS IN NUMERICAL LINEAR ALGEBRA}
\fancyhead[R]{\sffamily\bfseries\footnotesize\color{accent}SP-15}
\fancyfoot[L]{\parbox[b]{0.9\textwidth}{\sffamily\fontsize{8}{10}\selectfont\color{muted}Editorial ratings; a bounded search cannot certify that a problem remains unsolved.\\Literature check: 2026-09-12}}
\fancyfoot[R]{\sffamily\footnotesize\color{muted}\thepage}
\renewcommand{\headrulewidth}{0.3pt}
\renewcommand{\headrule}{\hbox to\headwidth{\color{accent}\leaders\hrule height \headrulewidth\hfill}}
\makeatletter
\renewcommand{\section}{\Needspace{5\baselineskip}\@startsection{section}{1}{0pt}{12pt plus 2pt minus 2pt}{4pt}{\sffamily\bfseries\large\color{ink}}}
\renewcommand{\subsection}{\Needspace{4\baselineskip}\@startsection{subsection}{2}{0pt}{9pt plus 1pt minus 1pt}{3pt}{\sffamily\bfseries\normalsize\color{ink}}}
\makeatother
\setcounter{secnumdepth}{0}
\begin{document}
{\sffamily\small\color{muted}Eigenvalues and inverse problems\par}
\vspace{3pt}
{\raggedright\hyphenpenalty=10000\exhyphenpenalty=10000\sffamily\bfseries\fontsize{21}{24}\selectfont\color{ink}Finitely
many unitary classes with prescribed shifted singular values\par}
\vspace{7pt}
{\sffamily\small\textbf{Difficulty:} challenging\quad\textbf{Importance:} interesting
to the community\par}
{\sffamily\small\bfseries\color{accent}Status: Solved\par}
\vspace{4pt}
{\color{accent}\hrule height 0.8pt}
\vspace{6pt}
\textbf{Topic:} Inverse determination from complete pseudospectral data

\textbf{Rating rationale:} Generic finiteness follows from invariant
theory, but the exceptional fibers resist that argument; the question
measures how much complete pseudospectral data can identify a nonnormal
matrix.

\subsection{Negative resolution - George
Stepaniants}\label{negative-resolution---george-stepaniants}

\textbf{Author:} George Stepaniants, Department of Computing and
Mathematical Sciences, California Institute of Technology, Pasadena,
California, USA.

\textbf{Solved negatively, 12 September 2026 (UTC).} The
\href{https://github.com/ajt60gaibb/OpenProblemsInNLA/blob/main/eigenvalues-and-inverse-problems/SP-15/solution.md}{Theorem
and Sections 1-4 of the complete proof} establish a continuous family of
complex \(9\times9\) matrices with the same singular values of every
complex scalar shift, with no two distinct family members unitarily
similar. Each matrix has nilpotent Jordan type \((3,3,3)\). Thus no
finite \(M_9\) exists, refuting the universal finiteness statement
below. This does not assert a classification of the remaining
dimensions.

\href{https://github.com/ajt60gaibb/OpenProblemsInNLA/blob/main/eigenvalues-and-inverse-problems/SP-15/solution.pdf}{Proof
PDF} ·
\href{https://github.com/ajt60gaibb/OpenProblemsInNLA/blob/main/eigenvalues-and-inverse-problems/SP-15/solution.tex}{Standalone
proof TeX} ·
\href{https://github.com/ajt60gaibb/OpenProblemsInNLA/blob/main/references/stepaniants-sp15-2026-09-12/verification/independent-review-aa01/review.md}{Independent
mathematical review} ·
\href{https://github.com/ajt60gaibb/OpenProblemsInNLA/blob/main/references/stepaniants-sp15-2026-09-12/README.md}{Submission
and eligibility record}.

The full proof passed a separate independent Codex-agent audit.
Substantial AI assistance is disclosed; this is informal automated
review, not external human peer review or Lean verification. The proof
uses an exact block determinant identity and a polynomial dimension
argument, with the standard constant-rank theorem. Its finite diagnostic
checks are not the proof. Fortier Bourque and Ransford retain credit for
the question and generic finiteness theorem, which is compatible with
this exceptional fiber.

\textbf{Retained historical material.} The original ratings, context,
exact target, references and dated pre-resolution status-search text are
preserved below. The earlier generic result is unchanged; the older
references to an unresolved question describe the pre-resolution record.

\newpage

\subsection{Context and notation}\label{context-and-notation}

For a complex \(N\times N\) matrix \(M\), write
\(s_1(M)\geq\cdots\geq s_N(M)\) for its singular values. Say
\(A\sim_{\rm sip} B\) if

\[
s_j(A-zI)=s_j(B-zI)
\qquad(z\in\mathbb C,\ 1\leq j\leq N).
\]

This is equality of \emph{super-identical pseudospectral} data: it
includes every singular value of every scalar shift.

\subsection{Problem statement}\label{problem-statement}

Is it true that for every \(N\geq1\) there is an integer \(M_N\geq1\)
such that every family \(A_1,\ldots,A_{M_N+1}\in\mathbb C^{N\times N}\)
with \(A_i\sim_{\rm sip}A_j\) for all \(i,j\) contains a pair satisfying

\[
A_j=U^*A_iU\quad\text{for some }i<j\text{ and }U^*U=I?
\]

Equivalently, is every such data fiber a union of at most \(M_N\)
unitary similarity classes, with a bound depending only on \(N\)?

\subsection{References}\label{references}

\begin{itemize}
\tightlist
\item
  Maxime Fortier Bourque and Thomas Ransford,
  \href{https://doi.org/10.1112/jlms/jdn085}{\emph{Super-identical
  pseudospectra}}, \emph{Journal of the London Mathematical Society}
  \textbf{79}(2), 511--528 (2009), \textbf{Theorem 1.4 and §6.2}. The
  theorem removes a closed measure-zero exceptional set; the question
  asks to remove this exception.
  \href{https://analyse.mat.ulaval.ca/abstracts/2008-03.pdf}{Author
  manuscript}.
\item
  Thomas Ransford,
  \href{https://analyse.mat.ulaval.ca/abstracts/2010-01.pdf}{\emph{Pseudospectra
  and matrix behaviour}}, \textbf{Theorem 5.4 and the following
  discussion, pp.~10--11}. This restates the finiteness theorem and
  explicitly leaves the empty-exceptional-set case open. It explains why
  replacing algebraicity with integrality fails to prove it.
\item
  G. Armentia, J. M. Gracia and F. E. Velasco,
  \href{https://doi.org/10.1016/j.laa.2011.01.014}{\emph{Identical
  pseudospectra of any geometric multiplicity}}, \emph{Linear Algebra
  and its Applications} \textbf{436}(6), 1683--1688 (2012). Their
  similarity theorem for super-identical pseudospectra gives ordinary
  similarity, which allows nonunitary changes of basis and does not
  answer the displayed question.
\item
  T. Ransford and N. Walsh,
  \href{https://arxiv.org/pdf/2109.14472v2}{\emph{A four-mean theorem
  and its application to pseudospectra}, arXiv:2109.14472v2}, 9 July
  2022, Theorems 1.3--1.4. Polynomial-norm comparisons and
  similarity-conditioning bounds do not establish finiteness of the
  unitary classes; the related norm question is retained in
  \href{https://github.com/ajt60gaibb/OpenProblemsInNLA/blob/main/matrix-functions-and-stability/MF-24/README.md}{MF-24}.
\end{itemize}

The original low-dimensional results give \(M_2=1\) and \(M_3=2\);
\(M_1=1\) is immediate. The unresolved question concerns exceptional
data fibers in higher dimensions, rather than generic matrices.

\subsection{Scope and status check}\label{scope-and-status-check}

The quantifiers above are an explicit restatement of the source question
about the exceptional set, with \(M_N\) counting classes instead of the
source's pigeonhole family size. The source is older: current openness
rests on bounded later-literature searches, not a recent explicit
reaffirmation. Searches through 2026-09-11 for super-identical
pseudospectra, finite unitary classes, the exceptional set and related
similarity results found no resolution. The 2012 ordinary-similarity
theorem and the cited 2022 norm estimates were checked for scope. They
do not remove the exceptional set.
\end{document}
